\subsection{Descent at the prime two}

We apply the coefficient statement in Lemma~\ref{lem:local-ring} to
the \(2\)-primary projection of the row. In the \(A=\Z_{(2)}\)
specialization used here, complex conjugation preserves the valuation
on \(A[\xi]\), so the norm determines
the valuation of the character value itself.

\begin{proposition}\label{prop:order}
If a real circulant Hadamard matrix has order \(n>1\), then
\(n=4u^2\) for an odd positive integer \(u\).
\end{proposition}

\begin{proof}
For the matrix under consideration, the row norm
\eqref{eq:hadamard-norm} has already given
\(n=2^{2s}u^2\) as in \eqref{eq:initial-order}. It remains to show
\(s=1\).

Project \(C_n\) onto \(C_{2^{2s}}\). The image \(T\) of \(h\) has
odd integer coefficients, each being a sum of \(u^2\) signs, and
\[
 TT^*=2^{2s}u^2.
\]
Suppose \(s\geq2\). We describe an operation starting with
\begin{equation}\label{eq:two-step}
 T\in\Z[C_{2^k}],\qquad TT^*=2^{2t}u^2,
 \qquad k\geq1,\quad t\geq2,
\end{equation}
where every coefficient of \(T\) is odd. Use
Lemma~\ref{lem:local-ring} with \(A=\Z_{(2)}\) and a primitive
\(2^k\)-th root \(\xi\). Complex conjugation preserves \(A[\xi]\)
and sends \(\xi-1\) to the associate \(-\xi^{-1}(\xi-1)\).
It therefore preserves the normalized valuation. Evaluating
\eqref{eq:two-step} gives
\[
 v(T(\xi))=t,\qquad T(\xi)\in 2^tA[\xi].
\]

Write \(T=\sum_{j=0}^{2^k-1}t_jX^j\). Its remainder modulo
\(\Phi_{2^k}(X)=1+X^{2^{k-1}}\) is
\[
 \sum_{j=0}^{2^{k-1}-1}(t_j-t_{j+2^{k-1}})X^j.
\]
The coefficient statement in Lemma~\ref{lem:local-ring}, with
\(\ell=t\), shows that every displayed difference belongs to
\(2^t\Z_{(2)}\). Each difference is an integer, so it is divisible
by \(2^t\) in \(\Z\).

Let \(q:\Z[C_{2^k}]\to\Z[C_{2^{k-1}}]\) be the group projection.
The coefficients of \(q(T)/2\) are
\[
 \frac{t_j+t_{j+2^{k-1}}}{2}.
\]
They are odd integers: the two odd summands are congruent modulo
\(4\), since \(t\geq2\). Thus \(T'=q(T)/2\) belongs to the smaller
integer group ring. Projection commutes with the involution, and division
of the projected norm by \(4\) gives
\[
 T'T'^*=2^{2(t-1)}u^2.
\]

Starting with \(k=2s\) and \(t=s\), perform this operation \(s-1\)
times. After \(j\) steps the indices are \(k=2s-j\) and \(t=s-j\);
every step starts with \(t\geq2\). The final element has
\(2^{s+1}\) odd integer coefficients and product with its star equal
to \(4u^2\). The identity coefficient of that product is the sum of the squares
of those coefficients. Each odd square is \(1\) modulo \(8\), and
\(2^{s+1}\) is divisible by \(8\). The sum is therefore \(0\)
modulo \(8\), whereas \(4u^2\) is \(4\) modulo \(8\).
This contradiction proves \(s=1\).
\end{proof}
